% docs/texte/05_theory/deconvolution.tex -- theory chapter "Deconvolution" (Interferosim)
% Version: 0.1.0 (2026-10-12)
% Notation as defined in texte/05_theory/notation.tex.
% ================================================================

\label{ch:deconvolution}
\index{deconvolution}\index{CLEAN}\index{Wiener filter}
Deconvolution estimates the true sky from the dirty image. Two
methods are implemented; which one is suitable depends on whether the
transfer function $\mathcal S$ has zeros.

\section{Wiener filter}

In the Fourier domain the measurement is
\begin{equation}
  D(u,v) = \mathcal S(u,v)\,\tilde I(u,v) + N(u,v)
\end{equation}
with the true image $\tilde I$, the transfer function $\mathcal S$ and
noise $N$. The linear filter $\mathcal W$ that minimises the mean square
error $E|\hat I - I|^2$ of $\hat I = \mathcal W D$ is the Wiener filter
\cite{wiener1949}
\begin{equation}
  \mathcal W = \frac{\mathcal S^*}{|\mathcal S|^2 + \Phi_n/\Phi_s}
  \label{eq:wiener}
\end{equation}
with the power spectra $\Phi_s$ of the signal and $\Phi_n$ of the
noise. For $\Phi_n\to0$ it becomes the inverse filter $1/\mathcal S$;
where the noise dominates it suppresses instead of amplifying.

\paragraph*{Filled aperture.} For a single dish of diameter $D$ the
transfer function is the autocorrelation of the circular aperture,
\begin{equation}
  \mathcal S(q) = \frac{2}{\pi}\left[\arccos\frac{q}{q_c}
     - \frac{q}{q_c}\sqrt{1-\frac{q^2}{q_c^2}}\right],
  \qquad q_c = \frac{D}{\lambda},
\end{equation}
non-zero everywhere below $q_c$. The Wiener filter compensates its
decrease and recovers resolution up to the diffraction limit
(\texttt{deconvolution.WienerFilter}).

\paragraph*{Sparse sampling.} Where $\mathcal S = 0$, Eq.~\eqref{eq:wiener}
gives $\mathcal W D = 0$ exactly. No linear filter can fill unmeasured
cells. For an interferometer, where $\mathcal S$ is either zero or the
number of samples $N_k$ of a cell, the filter reduces to
$\mathcal W_k = N_k/(N_k^2+K)$ with $K=\Phi_n/\Phi_s$: natural
weighting for $K\gg N_k^2$, uniform weighting for $K\ll N_k^2$. It is
therefore modelled as \texttt{imaging.WienerWeighting}.

\section{Högbom CLEAN}

CLEAN \cite{hogbom1974} assumes that the sky consists of a limited
number of point sources:
\begin{enumerate}
  \item Find the peak of the residual image (initially $I_D$).
  \item Subtract the dirty beam, shifted to the peak position and
        scaled by the loop gain $g$ (typically 0.05--0.2) times the
        peak value; store this as a CLEAN component.
  \item Repeat until the peak falls below a threshold or a maximum
        number of iterations is reached.
  \item Convolve the components with the \emph{clean beam} $B_C$ (a
        Gaussian fitted to the main lobe of $B_D$) and add the residual
        image.
\end{enumerate}
Because the component model has a Fourier transform defined
everywhere, CLEAN implicitly fills unmeasured $(u,v)$ cells. This
non-linear extrapolation is justified only if the sparsity assumption
holds; it is ideal for compact sources such as masers and poor for
extended emission, which is broken up into many components.

\section{Comparison}

\begin{center}
\begin{tabular}{@{}>{\raggedright\arraybackslash}p{3.4cm}
                >{\raggedright\arraybackslash}p{5.0cm}
                >{\raggedright\arraybackslash}p{5.0cm}@{}}
\toprule
 & Wiener filter & CLEAN \\
\midrule
Type & linear & non-linear (greedy) \\
Prior information & power spectrum of the sky & few positive point sources \\
Unmeasured cells & stay empty & filled from the model \\
Suitable for & filled aperture (dish) & sparse sampling (interferometer) \\
Fails for & gaps in the sampling & extended emission \\
\bottomrule
\end{tabular}
\end{center}
Notebook~A demonstrates both cases quantitatively.
